\chapter{Areas and Volumes}\label{ch:g7:areas}

Grade 6 measured rectangles, right triangles and boxes
(\cref{ch:g6:measure}). With a pair of scissors --- cut a piece here,
glue it there --- the same ideas now give the areas of parallelograms,
of \emph{all} triangles and of the disk, and the volumes of prisms and
cylinders.

\section{Area of a parallelogram}

\begin{theorem}[Parallelogram]\label{thm:g7:areas:parallelogram}
A parallelogram with a side of length $b$ (its \emph{base}) and
distance $h$ between that side and the opposite one (its
\emph{height}) has area
\[
A = b \times h .
\]
\end{theorem}

\begin{proof}[Proof by scissors]
Cut off the right-angled triangle sticking out on one side and glue it
on the other: the parallelogram becomes a $b \times h$ rectangle with
the same area.
\end{proof}

\begin{omfigure}
\begin{tikzpicture}[scale=0.82]
  \fill[omDef!14] (0,0) -- (4,0) -- (5.2,1.8) -- (1.2,1.8) -- cycle;
  \draw[thick] (0,0) -- (4,0) -- (5.2,1.8) -- (1.2,1.8) -- cycle;
  \fill[omThm!25] (0,0) -- (1.2,1.8) -- (1.2,0) -- cycle;
  \draw[omThm, thick] (0,0) -- (1.2,1.8) -- (1.2,0) -- cycle;
  \draw[dashed] (1.2,0) -- (1.2,1.8);
  \node[below, font=\small] at (2.0,0) {$b$};
  \draw[Stealth-Stealth] (6.0,0) -- (6.0,1.8)
    node[midway, right, font=\small] {$h$};
  \begin{scope}[xshift=8.2cm]
  \fill[omDef!14] (0,0) rectangle (4,1.8);
  \draw[thick] (0,0) rectangle (4,1.8);
  \fill[omThm!25] (2.8,0) -- (4,1.8) -- (4,0) -- cycle;
  \draw[omThm, thick] (2.8,0) -- (4,1.8) -- (4,0) -- cycle;
  \node[below, font=\small] at (2.0,0) {$b$};
  \end{scope}
\end{tikzpicture}

{\small The scissors proof: sliding the red triangle from the left end
to the right end turns the parallelogram into a rectangle --- same
area, $b \times h$.}
\end{omfigure}

\begin{remark}
The height is the \emph{perpendicular} distance between the two bases,
not the length of the slanted side! A very slanted parallelogram can
have long sides and a small area.
\end{remark}

\section{Area of a triangle}

\begin{theorem}[Triangle]\label{thm:g7:areas:triangle}
A triangle with base $b$ and corresponding height $h$ (the
perpendicular distance from the opposite vertex to the base line) has
area
\[
A = \frac{b \times h}{2} .
\]
\end{theorem}

\begin{proof}
Two copies of the triangle, one turned by a half-turn around the
midpoint of one side, fit together into a parallelogram with base $b$
and height $h$ (that is central symmetry at work,
\cref{prop:g7:central:preserved}). The triangle is half of it:
$\frac{bh}{2}$.
\end{proof}

\begin{omfigure}
\begin{tikzpicture}[scale=0.9]
  \fill[omDef!14] (0,0) -- (4,0) -- (1.4,2.1) -- cycle;
  \draw[thick] (0,0) -- (4,0) -- (1.4,2.1) -- cycle;
  \fill[omThm!12] (4,0) -- (1.4,2.1) -- (5.4,2.1) -- cycle;
  \draw[thick, omThm, dashed] (4,0) -- (5.4,2.1) -- (1.4,2.1);
  \draw[dashed] (1.4,0) -- (1.4,2.1);
  \draw[thin] (1.4,0) ++(-0.2,0) -- ++(0,0.2) -- ++(0.2,0);
  \node[below, font=\small] at (2,0) {$b$};
  \node[right, font=\small] at (1.42,1.0) {$h$};
\end{tikzpicture}

{\small A triangle and its half-turned copy (dashed) make a
parallelogram: the triangle's area is $\frac{b\times h}{2}$. Any of the
three sides can serve as the base --- with \emph{its own} height.}
\end{omfigure}

\begin{example}\label{ex:g7:areas:triangle}
A triangle has a base of $9$ cm and the corresponding height measures
$4$ cm:
\[
A = \frac{9 \times 4}{2} = \frac{36}{2} = 18 \text{ cm}^2 .
\]
For a \emph{right} triangle, the two legs are a base and its height:
we recover the formula of \cref{prop:g6:measure:formulas}.
\end{example}

\section{Area of a disk}

\begin{theorem}[Disk]\label{thm:g7:areas:disk}
A disk of radius $r$ has area
\[
A = \pi r^2 .
\]
\end{theorem}

\begin{proof}[Idea of proof]
Cut the disk into many thin equal slices and lay them head to tail:
they almost form a parallelogram of height $r$ whose base is half the
circumference, $\pi r$ (\cref{prop:g6:measure:circle}). Its area is
close to $\pi r \times r = \pi r^2$, and the approximation becomes
perfect as the slices get thinner. A complete proof needs integral
calculus, from the High School volume.
\end{proof}

\begin{omfigure}
\begin{tikzpicture}[scale=0.85]
  \foreach \a in {0,45,...,315}
    \fill[omDef!20, draw=omDef] (0,0) -- (\a:1.4)
      arc (\a:\a+45:1.4) -- cycle;
  \begin{scope}[xshift=4.3cm, yshift=-0.15cm]
  \foreach \i in {0,1,2,3} {
    \fill[omDef!20, draw=omDef]
      ({\i*1.1},0) -- ({\i*1.1+0.55},1.4) -- ({\i*1.1+1.1},0) -- cycle;
    \fill[omDef!20, draw=omDef]
      ({\i*1.1+0.55},1.4) -- ({\i*1.1+1.1},0) -- ({\i*1.1+1.65},1.4)
      -- cycle;
  }
  \draw[Stealth-Stealth] (5.2,0) -- (5.2,1.4)
    node[midway, right, font=\small] {$\approx r$};
  \node[below, font=\small] at (2.5,-0.1) {base $\approx \pi r$};
  \end{scope}
\end{tikzpicture}

{\small Slicing the disk and unrolling the slices: almost a
parallelogram of base $\pi r$ (half the border) and height $r$.}
\end{omfigure}

\begin{example}\label{ex:g7:areas:disk}
A round table top has radius $60$ cm. Its area is
$\pi \times 60^2 = 3600\pi \approx 11\,310$ cm$^2$, a bit more than
$1.1$ m$^2$. Watch the difference: the \emph{perimeter} $2\pi r$ uses
$r$, the \emph{area} $\pi r^2$ uses $r^2$.
\end{example}

\section{Prisms and cylinders}

\begin{definition}[Prism, cylinder]\label{def:g7:areas:prism}
A \emph{prism}\index{prism} is a solid with two identical parallel
polygonal faces (the \emph{bases}) joined by rectangles; a
\emph{cylinder}\index{cylinder} is the same with disks as bases. The
\emph{height} is the distance between the two bases.
\end{definition}

\begin{theorem}[Volumes]\label{thm:g7:areas:volume}
For a prism or a cylinder with base area $B$ and height $h$:
\[
V = B \times h .
\]
\end{theorem}
\admitted

\begin{example}\label{ex:g7:areas:volume}
A tent is a prism lying on its side: its bases are triangles of base
$2$ m and height $1.5$ m, and the tent is $3$ m long.
\begin{enumerate}
  \item Base area: $B = \frac{2 \times 1.5}{2} = 1.5$ m$^2$.
  \item Volume: $V = B \times h = 1.5 \times 3 = 4.5$ m$^3$.
\end{enumerate}
A tin can of radius $4$ cm and height $11$ cm:
$V = \pi \times 4^2 \times 11 = 176\pi \approx 553$ cm$^3$, roughly
half a liter.
\end{example}

\section{Exercises}

\begin{exercise}[$\star$]\label{exo:g7:areas:1}
Compute the area of a parallelogram with base $8$ cm and height
$5$ cm; of another with base $6.5$ m and height $4$ m.
\end{exercise}

\begin{exercise}[$\star$]\label{exo:g7:areas:2}
A parallelogram has sides $10$ cm and $6$ cm, and the height relative
to the $10$ cm base measures $4$ cm. Compute its area. Why is the
answer not $60$ cm$^2$?
\end{exercise}

\begin{exercise}[$\star$]\label{exo:g7:areas:3}
Compute the area of a triangle with base $12$ cm and height $7$ cm;
of a right triangle with legs $9$ and $6$; of a triangle with base
$5.5$ m and height $2$ m.
\end{exercise}

\begin{exercise}[$\star$]\label{exo:g7:areas:4}
Compute the area of a disk of radius $5$ cm, then of a half-disk of
radius $10$ cm (exact values with $\pi$, then rounded to the unit).
\end{exercise}

\begin{exercise}[$\star$]\label{exo:g7:areas:5}
Compute the perimeter \emph{and} the area of a disk of radius $3$ cm.
Which of the two answers is in cm and which in cm$^2$?
\end{exercise}

\begin{exercise}[$\star$]\label{exo:g7:areas:6}
Compute the volume of a prism with base area $24$ cm$^2$ and height
$10$ cm; of a cylinder of radius $3$ cm and height $8$ cm (exact,
then rounded).
\end{exercise}

\begin{exercise}[$\star$]\label{exo:g7:areas:7}
A triangle has area $28$ cm$^2$ and base $8$ cm. Find the
corresponding height, writing the equation first.
\end{exercise}

\begin{exercise}[$\star\star$]\label{exo:g7:areas:8}
Draw any triangle and measure carefully the three base--height pairs.
Compute $\frac{b \times h}{2}$ for each pair: the three results should
agree (up to measuring error). Why?
\end{exercise}

\begin{exercise}[$\star\star$]\label{exo:g7:areas:9}
A trapezoid-shaped garden can be split, by one diagonal, into two
triangles sharing the same height $h = 20$ m, with bases $30$ m and
$18$ m. Compute its total area. Can you guess a general formula for
the trapezoid?
\end{exercise}

\begin{exercise}[$\star\star$]\label{exo:g7:areas:10}
A cylindrical vase of radius $6$ cm contains water to a height of
$15$ cm. All the water is poured into an empty box
$18$ cm $\times$ $12$ cm at the base. What water height is reached in
the box? (Same volume, different base.)
\end{exercise}

\begin{exercise}[$\star\star$]\label{exo:g7:areas:11}
A pizza of diameter $30$ cm costs $9$ euros; one of diameter $40$ cm
costs $14$ euros. Compute the two areas and the price per $100$
cm$^2$. Which pizza is the better deal?
\end{exercise}

\begin{exercise}[$\star\star\star$]\label{exo:g7:areas:12}
The two legs of a right triangle measure $6$ cm and $8$ cm, and its
hypotenuse measures $10$ cm. Compute its area with the legs, then use
that area to find the height relative to the \emph{hypotenuse}.
\end{exercise}
